\section*{Abstract}
\addcontentsline{toc}{section}{Abstract}

% ~250 words. Write this LAST. Structure below is the order reviewers expect.
% NON-NEGOTIABLE: sentence 2 must state that cores are fixed at the historical
% baseline, so that the temporal signal is landscape-side. If a reader has to
% reach the Methods to learn that, they will feel misled. See CAU-002.

\noindent
\TODO{Context: cheetah range contraction, majority of range outside protected areas.} 

The southern African Cheetah (\textit{Acinonyx jubatus jubatus}), has had their range decreased every since their last genetic bottleneck around 12,000 - 18,000 years ago during the last mega-faunal extinction. This has only been exacerbated in the past decades due to both natural and anthropogenic reasons, leaving their current extant range majority outside protected areas. This leaves "barriers" like ranchers, roads and other human development in areas where cheetahs will be roaming. 

\TODO{Gap: existing regional work models habitat suitability, not observed change in
network permeability; the 2026 KAZA multispecies assessment excluded cheetah outright.}

Models looking at habitat suitability have been done in the past, but none in regards to observed change in the movement network permeability. One 2026 KAZA (Kavango-Zambezi Transfrontier Conservation Area) multi-species connectivity study excluded the cheetah outright in their study.

\TODO{What we did: held range cores fixed at the 2010--2016 confirmed-range baseline and
modelled \structconn\ across \nscen\ predeclared resistance scenarios at four snapshots
spanning \studyperiod.}
\TODO{Headline result: magnitude and geography of connectivity loss.}
\TODO{Second result: share of robust pinch points outside protected areas.}
\TODO{Third result: ranked monitoring cells and their stability across weightings.}
\TODO{Claim boundary sentence: these are modelled structural priorities, not
demonstrated cheetah movement.}

\vspace{1em}
\noindent\textbf{Keywords:} landscape connectivity; circuit theory; \textit{Acinonyx jubatus};
protected area coverage; monitoring prioritisation; southern Africa
